\subsection{Centered-Horner enclosure and selective Bernstein refinement}
\label{sec:p4-selective-hybrid-proof}
To obtain a meaningful competitor for the original
Bernstein-only enclosure, we add a midpoint-centered,
exact-rational interval Horner evaluator and a selective
two-stage certificate. Let the exact rational score-numerator
polynomial be $P(\lambda)=\sum_{j=0}^d a_j\lambda^j$
on $I=[\ell,r]\subseteq[0,1]$, with
$c=(\ell+r)/2$ and $h=(r-\ell)/2$. The translated
polynomial in $u=\lambda-c\in[-h,h]$ is
\[
P(c+u)=\sum_{k=0}^d b_k u^k,\qquad
b_k=\sum_{j=k}^d \binom{j}{k}a_jc^{j-k}.
\]
An interval Horner evaluation of these $b_k$
over $[-h,h]$ encloses every $P(\lambda)$ on $I$
by interval-arithmetic inclusion. Since the input
coefficients and endpoints are rational, our exact
\texttt{Fraction} arithmetic introduces no numerical
rounding at this step. Neither this fact nor
the power-basis coefficient identity certifies
the unverified floating EEG feature extractor.

\paragraph{Selective refinement algorithm.}
Evaluate centered Horner first. If its lower bound is
strictly positive or its upper bound strictly negative,
retain that sign without computing Bernstein
coefficients. If centered Horner is unresolved,
transform to the Bernstein basis on the \emph{same}
interval and accept a strict sign only if all its
coefficients share that strict sign. Otherwise
return unresolved. A polynomial whose coefficients
all vanish is a separately recorded exact tie.
Inclusion is mandatory at \emph{both}
stages; finite point samples cannot substitute
for either interval proof.

\paragraph{Certificate-union proposition.}
Suppose both the centered-Horner and Bernstein
procedures return a strict sign only when their exact
enclosures exclude zero. Then the selective method
is sound and certifies precisely the union of
the intervals certified by either procedure.
Indeed, a strict centered-Horner sign gives the
desired inclusion proof immediately. Otherwise the
algorithm invokes the Bernstein procedure, which
gives an independent inclusion proof if it succeeds.
If both procedures could certify opposite strict
signs on the same interval, the two sound
enclosures would force $P$ to be simultaneously
strictly positive and strictly negative on a
nonempty interval, which is impossible. No
unknown-sign interval is promoted to a certificate.
Consequently, for a fixed test collection $\mathcal I$,
\[
C_{\rm selective}(\mathcal I)
= C_{\rm centered}(\mathcal I)
\cup C_{\rm Bernstein}(\mathcal I),\qquad
|C_{\rm selective}|\geq
\max\{|C_{\rm centered}|,|C_{\rm Bernstein}|\}.
\]
This is a correctness property, \emph{not}
a theorem of runtime improvement. The hybrid
pays the centered-Horner cost for every interval
and the Bernstein cost on unresolved ones;
when centered Horner rarely resolves or
Bernstein is costly, this policy can be slower
than either stand-alone method.

The method and signed-root/tie tests are archived
in the public experiment source accompanying
workflow \texttt{38326786180}.
The strict mathematical scope is rational-input
polynomial sign determination. Certifying the
floating-point-generated EEG spectrum, tangent
features and original SPD eigensystem would
require further upstream error bounds.
